\documentclass[12pt]{article}
\usepackage[utf8]{inputenc}
\usepackage[T1]{fontenc}
\usepackage{hyperref}
\title{Mwani Language Resources for K-12 Curriculum Development}
\author{Compiled by Claude Sonnet 5.5 (claude-sonnet-5-5) for Andr\'as Kornai}
\date{}
\begin{document}
\maketitle

\section{Introduction}
Mwani (Kimwani; ISO 639-3 wmw) is a Bantu language spoken along the coast of Cabo Delgado, Mozambique, including the Quirimbas Islands. Search results describe about 60\% lexical similarity with Swahili but no mutual intelligibility. Kiwibo, the dialect of Ibo Island, is the prestige variety. Speaker estimates vary, from about 120,000 (Omniglot, 2006) to about 167,000. The search found no evidence of Mwani being taught in schools. Very little educational material turned up.

\section{Primary Educational Resources}
No primers, textbooks or school programmes for Mwani appeared in the search results. Any curriculum would have to be built from scratch, drawing on the linguistic description below.

\section{Linguistic and Reference Materials}
\begin{itemize}
\item Mwani: grammatical sketch (SIL Mo\c{c}ambique, 2010; hosted by Universit\"atsbibliothek Frankfurt). It follows the Kiwibo forms. \url{https://publikationen.ub.uni-frankfurt.de/opus4/frontdoor/index/index/year/2011/docId/22342}
\item Omniglot, Mwani (Kimwani): notes that Mwani was historically written in an Arabic-based script and that several Latin spellings are currently in use. No standard orthography was found. \url{https://www.omniglot.com/writing/mwani.htm}
\item Wikipedia, Mwani language: overview and classification. \url{https://en.wikipedia.org/wiki/Mwani_language}
\item Wikipedia, Mwani people: background on the community. \url{https://en.wikipedia.org/wiki/Mwani_people}
\end{itemize}
No dictionary was found in the search results.

\section{Religious and Community Resources}
A mission profile, Mwani of Mozambique, lists Bible portions, the Jesus film and some audio recordings, but no complete New Testament. It states that Scripture resources are still needed. \url{https://worldea.org/?p=14121}

\section{Digital and Audio Resources}
Omniglot gives a sample text. Nothing else digital was found beyond a number-counting page. \url{https://www.languagesandnumbers.com/how-to-count-in-mwani/en/wmw/}

\section{Curriculum Development Context}
Little published material exists. A curriculum would need to begin with orthography standardisation, using the SIL grammatical sketch as a starting point. It should involve local speakers on the Cabo Delgado coast. Mozambican schools generally teach in Portuguese; no Mwani-medium programme was found.

\end{document}